\documentclass[10pt,a4paper]{amsart}
\usepackage{amsmath,amssymb,amsthm,graphicx}
\usepackage{verbatim}
\newenvironment{reference}{\comment}{\endcomment}
%\newenvironment{reference}{}{}
\newenvironment{slogan}{\comment}{\endcomment}
\newenvironment{history}{\comment}{\endcomment}
% For commutative diagrams we use Xy-pic
\usepackage[all]{xy}

% We use 2cell for 2-commutative diagrams.
\xyoption{2cell}
% Theorem environments.
%
\theoremstyle{plain}
\newtheorem{theorem}[subsection]{Theorem}
\newtheorem{proposition}[subsection]{Proposition}
\newtheorem{lemma}[subsection]{Lemma}

\theoremstyle{definition}
\newtheorem{definition}[subsection]{Definition}
\newtheorem{example}[subsection]{Example}
\newtheorem{exercise}[subsection]{Exercise}
\newtheorem{situation}[subsection]{Situation}

\theoremstyle{remark}
\newtheorem{remark}[subsection]{Remark}
\newtheorem{remarks}[subsection]{Remarks}

\numberwithin{equation}{subsection}

% Macros
%
\def\lim{\mathop{\mathrm{lim}}\nolimits}
\def\colim{\mathop{\mathrm{colim}}\nolimits}
\def\Spec{\mathop{\mathrm{Spec}}}
\def\Hom{\mathop{\mathrm{Hom}}\nolimits}
\def\Ext{\mathop{\mathrm{Ext}}\nolimits}
\def\SheafHom{\mathop{\mathcal{H}\!\mathit{om}}\nolimits}
\def\SheafExt{\mathop{\mathcal{E}\!\mathit{xt}}\nolimits}
\def\Sch{\mathit{Sch}}
\def\Mor{\mathop{\mathrm{Mor}}\nolimits}
\def\Ob{\mathop{\mathrm{Ob}}\nolimits}
\def\Sh{\mathop{\mathit{Sh}}\nolimits}
\def\NL{\mathop{N\!L}\nolimits}
\def\CH{\mathop{\mathrm{CH}}\nolimits}
\def\proetale{{pro\text{-}\acute{e}tale}}
\def\etale{{\acute{e}tale}}
\def\QCoh{\mathit{QCoh}}
\def\Ker{\mathop{\mathrm{Ker}}}
\def\Im{\mathop{\mathrm{Im}}}
\def\Coker{\mathop{\mathrm{Coker}}}
\def\Coim{\mathop{\mathrm{Coim}}}

% Boxtimes
%
\DeclareMathSymbol{\boxtimes}{\mathbin}{AMSa}{"02}

%
% Macros for moduli stacks/spaces
%
\def\QCohstack{\mathcal{QC}\!\mathit{oh}}
\def\Cohstack{\mathcal{C}\!\mathit{oh}}
\def\Spacesstack{\mathcal{S}\!\mathit{paces}}
\def\Quotfunctor{\mathrm{Quot}}
\def\Hilbfunctor{\mathrm{Hilb}}
\def\Curvesstack{\mathcal{C}\!\mathit{urves}}
\def\Polarizedstack{\mathcal{P}\!\mathit{olarized}}
\def\Complexesstack{\mathcal{C}\!\mathit{omplexes}}
% \Pic is the operator that assigns to X its picard group, usage \Pic(X)
% \Picardstack_{X/B} denotes the Picard stack of X over B
% \Picardfunctor_{X/B} denotes the Picard functor of X over B
\def\Pic{\mathop{\mathrm{Pic}}\nolimits}
\def\Picardstack{\mathcal{P}\!\mathit{ic}}
\def\Picardfunctor{\mathrm{Pic}}
\def\Deformationcategory{\mathcal{D}\!\mathit{ef}}
\usepackage{cleveref}
\crefname{originallemma}{original lemma}{original lemmas}
\Crefname{originallemma}{Original lemma}{Original lemmas}
\makeatletter
\newlabel{derived-lemma-split}{{05QT}{source}{}{sourcepointer-1}{}}
\newlabel{derived-lemma-split@cref}{{[originallemma][0][]05QT}{[1][0][]source}{}{}{}}
\makeatother
\crefname{originallemma}{original lemma}{original lemmas}
\Crefname{originallemma}{Original lemma}{Original lemmas}
\makeatletter
\newlabel{divisors-lemma-immersion-smooth-into-smooth-regular-immersion}{{067U}{source}{}{sourcepointer-2}{}}
\newlabel{divisors-lemma-immersion-smooth-into-smooth-regular-immersion@cref}{{[originallemma][0][]067U}{[1][0][]source}{}{}{}}
\makeatother
\crefname{originallemma}{original lemma}{original lemmas}
\Crefname{originallemma}{Original lemma}{Original lemmas}
\makeatletter
\newlabel{lemma-distinguished-triangle-blowup}{{0FUV}{source}{}{sourcepointer-3}{}}
\newlabel{lemma-distinguished-triangle-blowup@cref}{{[originallemma][0][]0FUV}{[1][0][]source}{}{}{}}
\makeatother
\crefname{originallemma}{original lemma}{original lemmas}
\Crefname{originallemma}{Original lemma}{Original lemmas}
\makeatletter
\newlabel{lemma-gysin-via-log-complex}{{0FUQ}{source}{}{sourcepointer-4}{}}
\newlabel{lemma-gysin-via-log-complex@cref}{{[originallemma][0][]0FUQ}{[1][0][]source}{}{}{}}
\makeatother
\crefname{originallemma}{original lemma}{original lemmas}
\Crefname{originallemma}{Original lemma}{Original lemmas}
\makeatletter
\newlabel{more-morphisms-lemma-blowup}{{0FUT}{source}{}{sourcepointer-5}{}}
\newlabel{more-morphisms-lemma-blowup@cref}{{[originallemma][0][]0FUT}{[1][0][]source}{}{}{}}
\makeatother
\crefname{originallemma}{original lemma}{original lemmas}
\Crefname{originallemma}{Original lemma}{Original lemmas}
\makeatletter
\newlabel{morphisms-lemma-smooth-omega-finite-locally-free}{{02G1}{source}{}{sourcepointer-6}{}}
\newlabel{morphisms-lemma-smooth-omega-finite-locally-free@cref}{{[originallemma][0][]02G1}{[1][0][]source}{}{}{}}
\makeatother
\crefname{originalremark}{original remark}{original remarks}
\Crefname{originalremark}{Original remark}{Original remarks}
\makeatletter
\newlabel{remark-cup-product-as-a-map}{{0FU7}{source}{}{sourcepointer-7}{}}
\newlabel{remark-cup-product-as-a-map@cref}{{[originalremark][0][]0FU7}{[1][0][]source}{}{}{}}
\makeatother
\crefname{originalremark}{original remark}{original remarks}
\Crefname{originalremark}{Original remark}{Original remarks}
\makeatletter
\newlabel{remark-projective-space-bundle-formula}{{0FUN}{source}{}{sourcepointer-8}{}}
\newlabel{remark-projective-space-bundle-formula@cref}{{[originalremark][0][]0FUN}{[1][0][]source}{}{}{}}
\makeatother
\title{The de Rham complex of a blowup splits off the de Rham complex of the base}
\author{The Stacks Project Authors}
\date{stacks-project:derham.tex:local-snapshot-20261012; GNU FDL 1.2 or later}
\begin{document}
\maketitle
\noindent\textit{Editorial note.} All delivered mathematical text is the author's own, copied from the pinned Stacks Project chapter file derham.tex; no statement or proof step was rewritten, reordered or removed.\par\medskip
\noindent\textit{Editorial note.} Preamble provenance: the pinned chapter file opens with \textbackslash{}input\{preamble\} and carries no preamble text of its own, so the project's single Stacks-wide preamble staged at sources/stacks-derham/preamble.tex is excerpted verbatim as three spans, lines 8-12 (verbatim package and the project's own reference/slogan/history display environments), lines 14-18 (Xy-pic setup) and lines 154-217 (theorem declarations and mathematical macros). The wrapper class and the amsmath/amssymb/amsthm/graphicx base packages are supplied by the assembler. No macro was invented or redefined.\par\medskip
\noindent\textit{Editorial note.} Same-chapter pointers are located at their \textbackslash{}label lines inside the pinned chapter file; cross-chapter pointers are located in the project's own tag list staged in the snapshot (sources/stacks-chow/tags.txt), at the verbatim lines carrying their canonical tags.\par\medskip
\noindent\textit{Allowed standard prerequisites (editorial).}\begin{itemize}
\item The notation of the blowup and the distinguished triangle attached to it. More on Morphisms, label more-morphisms-lemma-blowup (cited at statement line 3189) and Same chapter derham.tex, label lemma-distinguished-triangle-blowup (invoked at proof line 3196). The first fixes the blowup diagram \$b : X' \textbackslash{}to X\$, \$i : Z \textbackslash{}to X\$, \$p : E \textbackslash{}to Z\$; the second supplies the triangle in the derived category of \$(X \textbackslash{}to S)\textasciicircum{}\{-1\}\textbackslash{}mathcal\{O\}\_S\$-modules from which the two summands below arise.
\item Splitting criterion in a triangulated category. Derived Categories, label derived-lemma-split (invoked at proof line 3203). It shows that a splitting of the map on the two summands whose image contains the summand \$i\_*\textbackslash{}Omega\textasciicircum{}\textbackslash{}bullet\_\{Z/S\}\$ yields a splitting of the first arrow of the triangle which vanishes on that summand.
\item Decomposition into components of fixed relative dimension, and regularity of the closed immersion. Morphisms, label morphisms-lemma-smooth-omega-finite-locally-free (invoked at proof line 3214) and Divisors, label divisors-lemma-immersion-smooth-into-smooth-regular-immersion (invoked at proof line 3221). They reduce to \$X\$ smooth of a fixed relative dimension \$N\$ and show that each \$i\_m : Z\_m \textbackslash{}to X\$ is a regular immersion of codimension \$N - m\$, so that the projective bundle formula applies to \$E\_m \textbackslash{}to Z\_m\$.
\item The projective space bundle formula and the description of its summands. Same chapter derham.tex, remark remark-projective-space-bundle-formula (invoked at proof lines 3235, 3251 and 3282). It supplies the isomorphism \$\textbackslash{}tilde\textbackslash{}xi\$ identifying \$Rp\_*\textbackslash{}Omega\textasciicircum{}\textbackslash{}bullet\_\{E/S\}\$ with a direct sum of shifts of \$\textbackslash{}Omega\textasciicircum{}\textbackslash{}bullet\_\{Z/S\}\$, together with the subcomplexes \$M\_m\$ and \$K\_m\$ satisfying \$M = K[-2]\$ and the commutativity of the diagram involving \$\textbackslash{}theta'\$.
\item The Gysin-type map of the blowup and its relation to the first Chern class of the exceptional line bundle. Same chapter derham.tex, label lemma-gysin-via-log-complex (invoked at proof line 3273) and remark remark-cup-product-as-a-map (invoked at proof line 3280). They provide the map \$\textbackslash{}delta : \textbackslash{}Omega\textasciicircum{}\textbackslash{}bullet\_\{E/S\}[-2] \textbackslash{}to \textbackslash{}Omega\textasciicircum{}\textbackslash{}bullet\_\{X'/S\}\$ whose composite with the restriction back to \$E\$ is \$\textbackslash{}theta'\$ for \$c\_1\textasciicircum{}\{dR\}(\textbackslash{}mathcal\{O\}\_E(1))\$, which is what makes the displayed chain a section.
\end{itemize}\medskip
\ifdefined\hypertarget\hypertarget{sourcepointer-1}{}\fi
\noindent\textit{Source reference.} Original lemma 05QT: sources/stacks-chow/tags.txt, source lines 7387--7387. Navigation pointer for \textbackslash{}ref\{derived-lemma-split\} used inside the counted span (line 3203); the argument that consumes it is retained in the span.\par\medskip
\ifdefined\hypertarget\hypertarget{sourcepointer-2}{}\fi
\noindent\textit{Source reference.} Original lemma 067U: sources/stacks-chow/tags.txt, source lines 7983--7983. Navigation pointer for \textbackslash{}ref\{divisors-lemma-immersion-smooth-into-smooth-regular-immersion\} used inside the counted span (line 3221); the argument that consumes it is retained in the span.\par\medskip
\ifdefined\hypertarget\hypertarget{sourcepointer-3}{}\fi
\noindent\textit{Source reference.} Original lemma 0FUV: sources/stacks-derham/pinned-original-source.tex, source lines 3132--3132. Navigation pointer for \textbackslash{}ref\{lemma-distinguished-triangle-blowup\} used inside the counted span (line 3196); the argument that consumes it is retained in the span.\par\medskip
\ifdefined\hypertarget\hypertarget{sourcepointer-4}{}\fi
\noindent\textit{Source reference.} Original lemma 0FUQ: sources/stacks-derham/pinned-original-source.tex, source lines 2560--2560. Navigation pointer for \textbackslash{}ref\{lemma-gysin-via-log-complex\} used inside the counted span (line 3273); the argument that consumes it is retained in the span.\par\medskip
\ifdefined\hypertarget\hypertarget{sourcepointer-5}{}\fi
\noindent\textit{Source reference.} Original lemma 0FUT: sources/stacks-chow/tags.txt, source lines 19788--19788. Navigation pointer for \textbackslash{}ref\{more-morphisms-lemma-blowup\} used inside the counted span (line 3189); the argument that consumes it is retained in the span.\par\medskip
\ifdefined\hypertarget\hypertarget{sourcepointer-6}{}\fi
\noindent\textit{Source reference.} Original lemma 02G1: sources/stacks-chow/tags.txt, source lines 3212--3212. Navigation pointer for \textbackslash{}ref\{morphisms-lemma-smooth-omega-finite-locally-free\} used inside the counted span (line 3214); the argument that consumes it is retained in the span.\par\medskip
\ifdefined\hypertarget\hypertarget{sourcepointer-7}{}\fi
\noindent\textit{Source reference.} Original remark 0FU7: sources/stacks-derham/pinned-original-source.tex, source lines 575--575. Navigation pointer for \textbackslash{}ref\{remark-cup-product-as-a-map\} used inside the counted span (line 3280); the argument that consumes it is retained in the span.\par\medskip
\ifdefined\hypertarget\hypertarget{sourcepointer-8}{}\fi
\noindent\textit{Source reference.} Original remark 0FUN: sources/stacks-derham/pinned-original-source.tex, source lines 2278--2278. Navigation pointer for \textbackslash{}ref\{remark-projective-space-bundle-formula\} used inside the counted span (line 3235); the argument that consumes it is retained in the span.\par\medskip
\begin{proposition}
\label{proposition-blowup-split}
With notation as in More on Morphisms, Lemma \ref{more-morphisms-lemma-blowup}
the map $\Omega^\bullet_{X/S} \to Rb_*\Omega^\bullet_{X'/S}$
has a splitting in $D(X, (X \to S)^{-1}\mathcal{O}_S)$.
\end{proposition}

\begin{proof}
Consider the triangle constructed in
Lemma \ref{lemma-distinguished-triangle-blowup}.
We claim that the map
$$
Rb_*(\Omega^\bullet_{X'/S}) \oplus i_*\Omega^\bullet_{Z/S} \to
i_*Rp_*(\Omega^\bullet_{E/S})
$$
has a splitting whose image contains the summand $i_*\Omega^\bullet_{Z/S}$.
By Derived Categories, Lemma \ref{derived-lemma-split} this will show that
the first arrow of the triangle has a splitting which vanishes on
the summand $i_*\Omega^\bullet_{Z/S}$ which proves the lemma.
We will prove the claim by decomposing $Rp_*\Omega^\bullet_{E/S}$
into a direct sum where the first piece corresponds to
$\Omega^\bullet_{Z/S}$ and the second piece can be lifted
through $Rb_*\Omega^\bullet_{X'/S}$.

\medskip\noindent
Proof of the claim. We may decompose $X$ into open and closed subschemes
having fixed relative dimension to $S$, see
Morphisms, Lemma \ref{morphisms-lemma-smooth-omega-finite-locally-free}.
Since the derived category $D(X, f^{-1}\mathcal{O})_S)$ correspondingly
decomposes as a product of categories, we may assume $X$ has
fixed relative dimension $N$ over $S$. We may decompose
$Z = \coprod Z_m$ into open and closed subschemes of relative
dimension $m \geq 0$ over $S$. The restriction $i_m : Z_m \to X$ of
$i$ to $Z_m$ is a regular immersion of codimension $N - m$, see Divisors, Lemma
\ref{divisors-lemma-immersion-smooth-into-smooth-regular-immersion}.
Let $E = \coprod E_m$ be the corresponding decomposition, i.e.,
we set $E_m = p^{-1}(Z_m)$. If $p_m : E_m \to Z_m$ denotes the
restriction of $p$ to $E_m$, then we have a canonical isomorphism
$$
\tilde \xi_m :
\bigoplus\nolimits_{t = 0, \ldots, N - m - 1}
\Omega^\bullet_{Z_m/S}[-2t]
\longrightarrow
Rp_{m, *}\Omega^\bullet_{E_m/S}
$$
in $D(Z_m, (Z_m \to S)^{-1}\mathcal{O}_S)$
where in degree $0$ we have the canonical map
$\Omega^\bullet_{Z_m/S} \to Rp_{m, *}\Omega^\bullet_{E_m/S}$.
See Remark \ref{remark-projective-space-bundle-formula}.
Thus we have an isomorphism
$$
\tilde \xi :
\bigoplus\nolimits_m
\bigoplus\nolimits_{t = 0, \ldots, N - m - 1}
\Omega^\bullet_{Z_m/S}[-2t]
\longrightarrow
Rp_*(\Omega^\bullet_{E/S})
$$
in $D(Z, (Z \to S)^{-1}\mathcal{O}_S)$
whose restriction to the summand
$\Omega^\bullet_{Z/S} = \bigoplus \Omega^\bullet_{Z_m/S}$ of the source
is the canonical map $\Omega^\bullet_{Z/S} \to Rp_*(\Omega^\bullet_{E/S})$.
Consider the subcomplexes $M_m$ and $K_m$ of the complex
$\bigoplus\nolimits_{t = 0, \ldots, N - m - 1} \Omega^\bullet_{Z_m/S}[-2t]$
introduced in Remark \ref{remark-projective-space-bundle-formula}.
We set
$$
M = \bigoplus M_m
\quad\text{and}\quad
K = \bigoplus K_m
$$
We have $M = K[-2]$ and by construction the map
$$
c_{E/Z} \oplus \tilde \xi|_M :
\Omega^\bullet_{Z/S} \oplus M
\longrightarrow
Rp_*(\Omega^\bullet_{E/S})
$$
is an isomorphism (see remark referenced above).

\medskip\noindent
Consider the map
$$
\delta : \Omega^\bullet_{E/S}[-2] \longrightarrow \Omega^\bullet_{X'/S}
$$
in $D(X', (X' \to S)^{-1}\mathcal{O}_S)$ of
Lemma \ref{lemma-gysin-via-log-complex}
with the property that the composition
$$
\Omega^\bullet_{E/S}[-2] \longrightarrow \Omega^\bullet_{X'/S}
\longrightarrow
\Omega^\bullet_{E/S}
$$
is the map $\theta'$ of Remark \ref{remark-cup-product-as-a-map} for
$c_1^{dR}(\mathcal{O}_{X'}(-E))|_E) = c_1^{dR}(\mathcal{O}_E(1))$.
The final assertion of Remark \ref{remark-projective-space-bundle-formula}
tells us that the diagram
$$
\xymatrix{
K[-2] \ar[d]_{(\tilde \xi|_K)[-2]} \ar[r]_{\text{id}} &
M \ar[d]^{\tilde x|_M} \\
Rp_*\Omega^\bullet_{E/S}[-2] \ar[r]^-{Rp_*\theta'} &
Rp_*\Omega^\bullet_{E/S}
}
$$
commutes. Thus we see that we can obtain the desired splitting of
the claim as the map
\begin{align*}
Rp_*(\Omega^\bullet_{E/S})
& \xrightarrow{(c_{E/Z} \oplus \tilde \xi|_M)^{-1}}
\Omega^\bullet_{Z/S} \oplus M \\
& \xrightarrow{\text{id} \oplus \text{id}^{-1}}
\Omega^\bullet_{Z/S} \oplus K[-2] \\
& \xrightarrow{\text{id} \oplus (\tilde \xi|_K)[-2]}
\Omega^\bullet_{Z/S} \oplus Rp_*\Omega^\bullet_{E/S}[-2] \\
& \xrightarrow{\text{id} \oplus Rb_*\delta}
\Omega^\bullet_{Z/S} \oplus Rb_*\Omega^\bullet_{X'/S}
\end{align*}
The relationship between $\theta'$ and $\delta$ stated above
together with the commutative diagram involving $\theta'$, $\tilde \xi|_K$,
and $\tilde \xi|_M$ above are exactly what's needed to
show that this is a section to the canonical map
$\Omega^\bullet_{Z/S} \oplus Rb_*(\Omega^\bullet_{X'/S}) \to
Rp_*(\Omega^\bullet_{E/S})$ and the proof of the claim is complete.
\end{proof}

\end{document}
